\chapter{跨讲推导训练}

先遮住推导。每题用一张纸写出适用规则、已给事实、缺项、结论和一项反例。
以下均为教学题设，题中认定的故意、风险、身份或数额条件视为给定输入。

\section{不作为、因果与罪过}

\paragraph{题设。}
甲是当班救生员，负责池内安全。乙失去浮力，甲见状，身边有能立即投送的救生圈。
甲为了使乙死亡，故意不投送。题设已确认在该时点投送可救下乙，乙最终因溺水死亡。
改题：甲发现乙时，已经没有现实可行的救助手段。哪一步首先改变？

\paragraph{推导。}
记 $D$ 为对乙有作为义务，$F$ 为有现实救助能力，$N$ 为未履行，$P$ 为履行可避免本次死亡，
$R$ 为该义务所防止的溺水风险在本次死亡中实现，$I$ 为杀人故意。
基于第五、七、八讲的教材分析，对此故意杀人的不作为模型，先分四步：
\begin{enumerate}
 \item “当班、负责池内安全、乙在池内遇险”落实对乙的救助义务 $D$。
 \item 可立即投送救生圈给出 $F$。故意不投送给出 $N$。故得到 $D\land F\land N$，进入不作为实行评价。
 \item 投送本可避免本次溺亡给出 $P$。死亡由义务所针对的溺水风险引起给出 $R$。两项合取支撑对这个死亡结果的归责。
 \item 希望乙死亡给出 $I$。在题设其他构成、违法与责任条件具备时，才得故意杀人既遂的评价。
\end{enumerate}
改题使 $F$ 为假，所以第2步即告失败。即使 $D,N,I$ 仍成立，也不能补回 $F$。
再变一次：只说救生圈在岸边，未给甲能否及时投送。此时把 $F$ 列为待证，结论仍须保留。
\textbf{记法：}义务、能做、没做、能救、归责、罪过。记住“能做”和“能救”是两道门。

\section{财产取得的三个近邻模型}

\paragraph{题设。}
货物由仓库管理员实际控制，金额等其他入罪条件给定。甲想非法取得货物，虚构已付款。
甲案一，管理员受骗并在其管理权限内将整箱货交给甲。甲案二，管理员没有交货，甲趁其不备自行取走。
甲案三，所有人先正常交甲保管，甲后来才起意据为己有，经催还仍拒不退还。

\paragraph{推导。}
记 $E$ 为欺骗造成错误，$V$ 为基于错误实施有权或有地位的财产处分，$T$ 为被告自行破坏他人占有并建立控制，
$K$ 为此前已经代为保管，$A$ 为此后非法据为己有。
原书对处分意识及三角关系的细化规则在第十九讲，以下用其一般分析框架。
\[
\begin{array}{c|ccccc|l}
 \text{案}&E&V&T&K&A&\text{首先命中的路径}\\\hline
 \text{一}&1&1&0&0&0&\text{欺骗型交付}\\
 \text{二}&*&0&1&0&0&\text{自行转移占有}\\
 \text{三}&0&0&0&1&1&\text{已持有后转为非法占有}
\end{array}
\]
表中 $*$ 表示这一项对第二案的区分不必先定值。案一由错误到交付，再由交付到取得与损失，满足诈骗的核心因果链。
案二没有处分环节，虚构付款本身不能补出 $V$，应审查盗窃路径。
案三取财前已基于保管合法控制财物，判断侵占还须加入题设给出的拒不退还及法定门槛。
三个结论均是在相应主观、责任等条件已具备的前提下得出。
\textbf{记法：}画取财前后两帧，问“原来谁占有，谁改变占有，交付基于什么”。

\section{共同故意、实行过限和转化型抢劫}

\paragraph{题设。}
甲、乙共同商定盗窃。乙按分工在外围望风，甲已取走财物并受到现场追赶。
甲为抗拒抓捕当场以暴力威胁追赶人。乙没有在事前或事中接受、参与该暴力升级。
假定先行行为及暴力强度满足《刑法》第269条及适用司法解释的要求，其他责任条件已具备。

\paragraph{推导。}
用 $T(p)$ 表示先行盗窃的参与，用 $A(p)$ 表示当场，用 $M(p)$ 表示窝藏赃物、抗拒抓捕或者毁灭罪证的法定目的，
用 $V(p)$ 表示实施暴力或以暴力相威胁。对实行者，第269条的核心路径为
\[
 T(p)\land A(p)\land M(p)\land V(p)
 \ \Rightarrow\ \mathrm{TreatAsRobbery}(p).
\]
甲满足四个条件，得到按抢劫定罪处罚的分支。乙满足共同盗窃的分工与故意，
但追加暴力不在乙已确认的共同范围，不能从“共同盗窃”推出“共同抢劫”。对乙应按参与盗窃的事实单独判断。
若乙当场得知并接受暴力方案，主动拦挡追赶者配合甲，第一变化发生于共同故意和行为贡献。
此时应重新审查乙对转化行为的共同参与，仍要核对时间和作用。
\textbf{记法：}先画“共同范围”，再画“追加行为”。箭头越界时，逐人找接受和贡献。

\section{数罪并罚：漏罪与新罪改变运算顺序}

\paragraph{题设。}
原判有期徒刑十年，已执行六年。第二项单独判有期徒刑八年。均无其他刑种或特别限制。
两种情形分别是：第二项属于原判宣告前实施的漏罪，或者属于原判后、执行完毕前实施的新罪。
求剩余执行刑期的许可区间。

\paragraph{推导。}
对两个有期徒刑，令 $C(x,y)$ 表示第69条许可的决定执行区间：
\[
C(x,y)=
 \left[\max(x,y),\min\{x+y,U(x+y)\}\right],
\quad
 U(s)=\begin{cases}20,&s<35,\\25,&s\ge35.\end{cases}
\]
式中刑期均用年作单位，边界可按月进一步比较。本题中：
\begin{align*}
 \text{漏罪（第70条）}\quad
 C(10,8)&=[10,18],&
 [10,18]-6&=[4,12];\\
 \text{新罪（第71条）}\quad
 10-6&=4,&
 C(4,8)&=[8,12].
\end{align*}
第一行是新判决决定的刑期内计入已执行部分。第二行是将前罪未执行部分与后罪刑期合并。
即使上界相同，下界也随顺序改变。另设三罪分别判处十五年、十年和十年，刑期总和恰好三十五年，上限函数进入二十五年分支，合并区间为十五至二十五年。
\textbf{记法：}先问行为在原判前还是后。漏罪先并后减，新罪先减后并。

\section{公款行贿、渎职及并罚}

\paragraph{题设。}
国家工作人员甲利用职务便利，用公共财物向他人行贿。题设明确无个人占有目的，
并分别给出行贿罪和相应渎职犯罪的全部构成、主观责任及其他条件。
改题：只知有公共财物给付，渎职要求的严重结果及归责均未证明。

\paragraph{推导。}
令 $N,U,P$ 分别表示国家工作人员、利用职务便利和使用公共财物，$A$ 表示个人占有目的，
$B,G$ 分别表示行贿罪和渎职犯罪已经成立。法释〔2026〕6号第20条的特别处罚规则可记为
\[
 N\land U\land P\land\neg A\land B\land G
 \ \Rightarrow\ \mathrm{Combine}(B,G).
\]
先依两条构成路径分别得到 $B,G$，再合取其余已给事实，适用并罚规则。
改题中 $G$ 未能推出，故不能经此规则得到并罚，也不能只由“渎职未证实”否定行贿罪。
若改为先个人占有公款再用所得行贿，$\neg A$ 的前提失效，应另查贪污及行贿各自的成立条件和罪数关系。
\textbf{记法：}公款走向、占有目的、两罪各证、最后并罚。
